% =====================================================================
%  FICHE 8 — THÈME 3 — NIVEAU DIFFICILE — CORRIGÉ
% =====================================================================
\documentclass[11pt,a4paper]{article}
\usepackage[utf8]{inputenc}
\usepackage[T1]{fontenc}
\usepackage{lmodern}
\usepackage[french]{babel}
\usepackage{amsmath,amssymb,mathtools}
\usepackage{icomma}
\usepackage{siunitx}
\sisetup{output-decimal-marker={,},per-mode=power,group-separator={\,},group-minimum-digits=5}
\usepackage[version=4]{mhchem}
\usepackage{xcolor}
\usepackage{tikz}
\usepackage[most]{tcolorbox}
\usepackage{enumitem}
\usepackage{array}
\usepackage{fancyhdr}
\usepackage[hidelinks]{hyperref}
\usetikzlibrary{arrows.meta,positioning,calc}
\usepackage[a4paper,margin=2.1cm,headheight=26pt,headsep=10pt]{geometry}
\usepackage{titlesec}

\definecolor{primary}{RGB}{150,98,162}
\definecolor{primarydark}{RGB}{92,52,110}
\definecolor{excol}{RGB}{0,135,90}

\tcbset{boxbase/.style={enhanced,breakable,boxrule=0.9pt,arc=2.5pt,left=3mm,right=3mm,top=2.3mm,bottom=2.3mm,fonttitle=\bfseries,coltitle=white,toptitle=0.6mm,bottomtitle=0.6mm}}
\newtcolorbox[auto counter]{correction}[1][]{boxbase,colback=excol!4,colframe=excol,title={\textsc{Corrigé}~\thetcbcounter\ \ifx\relax#1\relax\relax\else\textmd{--- \itshape #1}\fi}}

\newcommand{\rep}[1]{\textcolor{primarydark}{\textbf{#1}}}

\pagestyle{fancy}
\fancyhf{}
\fancyhead[L]{\small\textcolor{primarydark}{\textbf{Fiche 8} — Thème 3 : Calorimétrie}}
\fancyhead[R]{\small\textcolor{primarydark}{\textsc{Difficile — corrigé}}}
\fancyfoot[C]{\small\thepage}
\renewcommand{\headrulewidth}{0.8pt}
\renewcommand{\headrule}{\hbox to\headwidth{\color{primary}\leaders\hrule height \headrulewidth\hfill}}
\setlength{\parindent}{0pt}\setlength{\parskip}{3pt}

\begin{document}

\begin{center}
\begin{tikzpicture}
\node[rectangle,rounded corners=7pt,fill=excol,text=white,minimum width=16cm,
      minimum height=1.1cm,align=center,inner sep=6pt]
  {{\large\bfseries Thème 3 — Calorimétrie}\\[1pt]{\small Niveau \textsc{Difficile} — corrigé détaillé}};
\end{tikzpicture}
\end{center}

\begin{correction}[de la glace à la vapeur — bilan complet]
$m=\SI{0.050}{\kilogram}$.
\begin{enumerate}[label=\alph*),leftmargin=2em,topsep=2pt,itemsep=2pt]
  \item \textcircled{1} échauffement glace $-40\to0$ : $m c_{\text{glace}}\Delta T$ ; \textcircled{2} fusion : $mL_{\text{fus}}$ ;
        \textcircled{3} échauffement eau $0\to100$ : $m c_{\text{eau}}\Delta T$ ; \textcircled{4} vaporisation : $mL_{\text{vap}}$ ;
        \textcircled{5} échauffement vapeur $100\to120$ : $m c_{\text{vap}}\Delta T$.
  \item $Q_1 = 0{,}050\times2{,}060\times40 = \SI{4.12}{\kilo\joule}$ ;\;
        $Q_2 = 0{,}050\times334 = \SI{16.7}{\kilo\joule}$ ;\;
        $Q_3 = 0{,}050\times4{,}185\times100 = \SI{20.925}{\kilo\joule}$ ;\\
        $Q_4 = 0{,}050\times2257 = \SI{112.85}{\kilo\joule}$ ;\;
        $Q_5 = 0{,}050\times1{,}85\times20 = \SI{1.85}{\kilo\joule}$.
  \item $Q = 4{,}12+16{,}7+20{,}925+112{,}85+1{,}85 = \rep{\SI{156.445}{\kilo\joule}}$.
  \item L'étape dominante est la \rep{vaporisation} ($Q_4=\SI{112.85}{\kilo\joule}$), soit
        $\dfrac{112{,}85}{156{,}445}\approx \rep{72\,\%}$ du total : $L_{\text{vap}}$ est bien plus grand
        que les autres termes.
\end{enumerate}
\end{correction}

\begin{correction}[barre de fer froide dans l'eau — congélation partielle]
\begin{enumerate}[label=\alph*),leftmargin=2em,topsep=2pt,itemsep=2pt]
  \item $Q_{\text{fer}} = m_{\text{fer}}c_{\text{fer}}\Delta T = 0{,}500\times0{,}5\times\big(0-(-34)\big)
        = 0{,}500\times0{,}5\times34 = \rep{\SI{8.5}{\kilo\joule}}$ (énergie à \emph{fournir} au fer).
  \item $|Q_{\text{eau}}| = m_{\text{eau}}c_{\text{eau}}|\Delta T| = 0{,}200\times4\times5 = \rep{\SI{4}{\kilo\joule}}$ (cédée en refroidissant).
  \item Le fer réclame \SI{8.5}{\kilo\joule} mais l'eau ne peut en céder que \SI{4}{\kilo\joule} en
        atteignant \SI{0}{\celsius}. Il \rep{manque $8{,}5-4=\SI{4.5}{\kilo\joule}$}, qui ne peuvent
        provenir que de la \rep{congélation} d'une partie de l'eau (la solidification libère la chaleur latente).
  \item Si \emph{toute} l'eau gelait, elle libérerait $0{,}200\times334 = \SI{66.8}{\kilo\joule}$, bien
        plus que les \SI{4.5}{\kilo\joule} requis : donc \rep{seule une partie gèle}. Bilan :
        $8{,}5 = 4 + m\times334 \Rightarrow m = \dfrac{4{,}5}{334} = \SI{0.01347}{\kilogram} \approx \rep{\SI{13.5}{\gram}}$.
  \item Température finale : \rep{\SI{0}{\celsius}} (coexistence glace–eau). Masse restée liquide :
        $200 - 13{,}5 = \rep{\SI{186.5}{\gram}}$.
\end{enumerate}
\end{correction}

\begin{correction}[dissolution avec un calorimètre qui s'échauffe]
\begin{enumerate}[label=\alph*),leftmargin=2em,topsep=2pt,itemsep=2pt]
  \item $\Delta T = 23{,}5 - 21{,}0 = \rep{\SI{2.5}{\kelvin}}$.
  \item $Q_1 = m_e c_e \Delta T = 0{,}150\times4{,}185\times2{,}5 = \rep{\SI{1.569}{\kilo\joule}}$.
  \item Seule la moitié du verre s'échauffe : $Q_2 = \dfrac{m_v}{2}\,c_v\,\Delta T
        = 0{,}100\times0{,}84\times2{,}5 = \rep{\SI{0.21}{\kilo\joule}}$.
  \item Chaleur totale captée : $Q = Q_1 + Q_2 = 1{,}569 + 0{,}21 = \SI{1.779}{\kilo\joule}$. La dissolution
        réchauffe le milieu : elle est \emph{exothermique}, elle a donc \emph{cédé} cette énergie. Par mole :
        \[ \Delta H = \frac{-\,1{,}779}{0{,}10} = \rep{\SI{-17.79}{\kilo\joule\per\mol}} \approx \SI{-17.8}{\kilo\joule\per\mol}. \]
\end{enumerate}
\end{correction}

\end{document}
